\chapter{Estudos Complementares}
\label{sec:complementares}

O Capítulo~\ref{sec:results} reportou a comparação principal na suíte sintética. Este capítulo reúne as famílias de evidência que examinam a QP1 fora dessa suíte e as outras quatro questões, uma seção por questão: a transferência para problemas de engenharia (QP1), a calibração e a ablação (QP2), a comparação de cada acoplamento com o seu hospedeiro (QP3), a decisão de quando ativar o operador (QP4) e o posicionamento contra os algoritmos comparados (QP5). As respostas às cinco questões estão no Capítulo~\ref{sec:discussion_conclusions}.

Cada família tem coorte e correção próprias, declaradas na Tabela~\ref{tab:familias_evidencia}, e as regras da Seção~\ref{sec:familias} valem em todas as seções; onde há sobreposição de coortes, ela é declarada na abertura da seção. As tabelas das famílias de engenharia, de calibração e ablação e do controlador são regeneradas a partir dos artefatos processados.

\section{Transferência para problemas de engenharia}
\label{sec:engenharia}

Esta seção verifica a QP1 fora da suíte sintética: o ganho sobre o hospedeiro se mantém em problemas com restrições? A evidência é de apoio externo e tem força inferior à da comparação principal: são três problemas, e não uma suíte.

Os critérios de inclusão relatados no manuscrito de origem são: dois ou três objetivos, o mesmo orçamento da suíte sintética, soluções viáveis não dominadas para o IVF/SPEA2 e disponibilidade de IGD e HV com cobertura de execuções comuns~\cite{zambrano2026ivfspea2}. Quatro candidatos não atenderam ao critério de viabilidade: no RWMOP13, no RWMOP20, no RWMOP24 e no RWMOP29, o IVF/SPEA2 não produziu solução viável em nenhuma de cinco execuções. O RWMOP9 foi mantido como problema de continuidade, segundo o manuscrito, e o RWMOP8 foi mantido embora a triagem já o mostrasse desfavorável no posicionamento entre algoritmos. A retenção de um caso adverso limita o viés de seleção, sem eliminá-lo.

\input{generated/tab_engenharia_wtl}

\begin{figure}[!htb]
    \centering
    \includegraphics[width=\textwidth]{generated/figures/fig_engenharia.pdf}
    \caption[Perfis de IGD e HV nos três problemas de engenharia]{Perfis de IGD e HV nos três problemas de engenharia.}
    \label{fig:engenharia}
\end{figure}

A Tabela~\ref{tab:engenharia_wtl} separa duas leituras que uma apresentação agregada confundiria. Contra o próprio hospedeiro, que é a comparação da QP1, o IVF/SPEA2 registra uma vitória detectável no RWMOP9, em IGD e em HV, e nenhuma diferença detectável nos outros dois problemas. A vitória do RWMOP9 sobrevive à correção de Holm entre os três problemas, mas tem magnitude desigual nos dois indicadores: a IGD mediana melhora $1{,}17\%$, e o HV mediano, $0{,}01\%$.

Contra o conjunto dos comparadores, a leitura é outra, e é a que o manuscrito de origem reporta~\cite{zambrano2026ivfspea2}. No RWMOP9 o IVF/SPEA2 vence todos os oito em IGD, mas em HV vence cinco e perde três: proximidade e cobertura não concordam nesse problema, como ilustra a Figura~\ref{fig:engenharia}. O RWMOP21 apresenta perfil equilibrado, com três vitórias e três derrotas nos dois indicadores. O RWMOP8 é adverso, com uma vitória e cinco derrotas entre os sete comparadores que têm execuções válidas. Esse posicionamento é misto, mas responde a outra pergunta: como o acoplamento se situa entre algoritmos distintos, e não se ele melhora o seu hospedeiro.

Na Figura~\ref{fig:engenharia}, cada ponto é a mediana de um algoritmo, as barras vão do primeiro ao terceiro quartil, e o IVF/SPEA2 aparece destacado. O rótulo à direita de cada linha indica o número de execuções com métricas válidas, de 60 previstas por algoritmo. No RWMOP8, a cobertura é heterogênea, o que restringe a comparação com o SPEA2+SDE, que tem 18 execuções válidas, e com o MOEA/D, que não tem nenhuma.

Nos três problemas avaliados, foi detectada melhora em um e nenhuma piora detectável nos outros dois; nesses dois, a ausência de diferença detectável não exclui prejuízo ao SPEA2. Entre os demais algoritmos, a sua posição varia de problema para problema.

\section{Calibração e ablação}
\label{sec:calibracao_ablacao}

Esta seção reúne a evidência da QP2: as diferenças de desempenho entre as configurações comparadas na calibração e entre as variantes de acoplamento comparadas na ablação. A calibração compara a configuração promovida com a melhor configuração da fase inicial e com a configuração inicial não ajustada; a ablação compara a variante com as duas decisões de acoplamento com o SPEA2 canônico e com a formulação sem essas decisões.

A Tabela~\ref{tab:tuning_c26} apresenta as comparações par a par sobre o subconjunto de calibração descrito na Seção~\ref{sec:calibracao}.

\input{generated/tab_tuning_c26}

A Figura~\ref{fig:calibracao} resume a varredura das três fases sobre as mesmas 12 configurações problema--objetivo, com 30 execuções por configuração. Nos painéis~(a) e~(c), cada célula é o posto médio combinado de IGD e HV de uma configuração entre os candidatos da fase; no painel~(b), cada barra é o mesmo posto para um dos cinco perfis de manipulação genética. Postos menores, em cores mais claras, são melhores, e a estrela marca a configuração vencedora de cada fase.

\begin{figure}[!htb]
    \centering
    \includegraphics[width=\textwidth]{generated/figures/fig_calibracao.pdf}
    \caption[Panorama da varredura de calibração em três fases]{Posto médio combinado de IGD e HV nas três fases da calibração.}
    \label{fig:calibracao}
\end{figure}

A configuração promovida não apresenta diferença detectável em relação à melhor configuração da fase inicial em nenhuma das 12 instâncias, em IGD, e vence em apenas uma delas em HV. As duas configurações apresentam quatro vitórias corrigidas em IGD diante da configuração inicial não ajustada, sem derrota detectável. Esse comparador histórico difere das configurações calibradas em mais do que os parâmetros: corresponde à formulação anterior do operador, sem as duas decisões de acoplamento, reúne 60 execuções por instância, contra 30 das configurações calibradas, e seu orçamento não consta do artefato desta comparação (Tabela~\ref{tab:tuning_c26}). A Figura~\ref{fig:calibracao} mostra postos próximos na vizinhança da configuração promovida, mas nem a figura nem a ausência de diferença detectável demonstram equivalência entre as configurações.

A ablação fatorial combinou as duas decisões de acoplamento com outras modificações candidatas sobre o subconjunto de calibração. O teste de Friedman~\cite{friedman1937use} não detectou diferença entre as 16 combinações ($p = 0{,}485$); a de melhor posto médio, a que adota as duas decisões e nenhuma outra modificação, foi levada às 51 instâncias sintéticas. A Tabela~\ref{tab:ablacao} compara essa variante com o SPEA2 canônico e com a formulação sem as decisões descritas nas Seções~\ref{sec:pai_dissimilar} e~\ref{sec:critério_coletivo}.

\input{generated/tab_ablacao}

O resultado contra o SPEA2 canônico é favorável: 25 vitórias em IGD e 28 em HV, com uma e duas derrotas. Contra a formulação sem as duas decisões, a variante registra 50 resultados sem diferença detectável e uma derrota em cada indicador, o que não demonstra benefício das decisões nem equivalência entre as formulações. Como registra a nota da Tabela~\ref{tab:ablacao}, a ablação foi executada na configuração inicial não ajustada, e não na configuração promovida que o Capítulo~\ref{sec:results} avalia, e o braço sem as duas decisões provém de outra campanha, sem pareamento por semente. A Seção~\ref{sec:resposta_qp2} reúne as comparações da calibração e da ablação e discute o que elas não permitem atribuir.

\section{Comparação dos acoplamentos com seus hospedeiros}
\label{sec:hospedeiros}

Esta seção reúne a evidência da QP3: como varia o desempenho relativo dos acoplamentos IVF/SPEA2, IVF/NSGA-II e IVF/NSGA-III diante dos respectivos hospedeiros. Cada um desses acoplamentos publicados é testado contra o seu próprio hospedeiro, sob protocolo comum de 30 execuções e correção de Benjamini--Hochberg por acoplamento e métrica; a comparação entre os três se faz sobre os resultados desses testes, sem teste que os compare diretamente entre si.

Três aspectos do desenho condicionam a leitura dos números. O primeiro é a sobreposição: as execuções 3001--3030 e 1--30 do par IVF/SPEA2 são um subconjunto da coorte da comparação principal do Capítulo~\ref{sec:results}, e só as execuções 4001--4030 e 5001--5030 acrescentam observações novas (Seção~\ref{sec:validade_conclusao}). O segundo é a unidade de comparação: cada coluna é o par formado por um hospedeiro e pela implementação, neste projeto, da realização do operador que os seus autores publicaram para ele, com diferenças de seleção do doador primário, de operador de recombinação, de critério de continuação e de regra de ativação. Diferem também a configuração e o modo como ela foi escolhida. A do IVF/SPEA2 é a calibrada em 12 das 51 instâncias desta suíte (Seção~\ref{sec:calibracao}). Os pipelines NSGA usaram configurações fixas, sem reajuste por instância: taxa de $0{,}5$, coleta de 7\% da população e até cinco ciclos no IVF/NSGA-II; taxa e coleta de 10\% e até cinco ciclos no IVF/NSGA-III. Esses valores diferem em parte dos protocolos de origem e não têm registro de ajuste equivalente ao do IVF/SPEA2 nesta suíte. A comparação é, portanto, entre pipelines assim configurados e não isola causalmente o efeito do hospedeiro. O terceiro é o pareamento: os pares NSGA compartilham a faixa de execuções com o hospedeiro e são pareados execução a execução, ao passo que, no par IVF/SPEA2, as faixas não coincidem e as execuções são alinhadas por ordem; a nota da Tabela~\ref{tab:hosts_wtl} registra o efeito de dispensar esse alinhamento.

\input{generated/tab_hosts_wtl}

A ordenação descritiva entre os três acoplamentos é a mesma nos dois indicadores, e a Figura~\ref{fig:hospedeiros_a12} exibe os resultados instância a instância. O IVF/SPEA2 melhora o seu hospedeiro em 33 das 51 instâncias em IGD e em 37 em HV. O IVF/NSGA-III melhora em 22 e 21, sem derrota em IGD, mas com uma proporção alta de empates. O IVF/NSGA-II registra uma vitória, 48 comparações sem diferença detectável e duas derrotas em IGD.

Na Figura~\ref{fig:hospedeiros_a12}, cada ponto é o $A_{12}$ de Vargha--Delaney do IGD numa das 51 instâncias, calculado sobre as 30 execuções de cada algoritmo e orientado de modo que valores acima de $0{,}5$ favoreçam o acoplamento; o traço horizontal é a mediana de cada coluna.

\begin{figure}[!htb]
    \centering
    \includegraphics[width=0.92\textwidth]{generated/figures/fig_hospedeiros_a12.pdf}
    \caption[Tamanho de efeito do IGD entre cada acoplamento IVF e o seu hospedeiro]{$A_{12}$ orientado do IGD, instância a instância, entre cada acoplamento IVF e o seu hospedeiro.}
    \label{fig:hospedeiros_a12}
\end{figure}

A Tabela~\ref{tab:hosts_geometria} estratifica essas contagens pela geometria da fronteira de Pareto.

\input{generated/tab_hosts_geometria}

A estratificação é consistente entre os três hospedeiros no sentido do efeito: todos rendem mais em fronteiras regulares do que em irregulares. Para o IVF/SPEA2, o tamanho de efeito mediano cai de $0{,}781$ para $0{,}642$ em IGD quando se passa das 40 instâncias regulares às 11 irregulares. A magnitude da queda, porém, não acompanha a magnitude do ganho: a maior é a do IVF/NSGA-III, de $0{,}687$ para $0{,}529$ em IGD, e não a do IVF/SPEA2, que é o acoplamento que mais ganha. A menor é a do IVF/NSGA-II, que quase não ganha em nenhum dos dois grupos.

Essa estratificação é descritiva. Nesta família, testou-se apenas a interação entre hospedeiro e geometria, que não atingiu significância, e onze instâncias irregulares, que correspondem a nove funções (Seção~\ref{sec:suites}), sustentam no máximo um gradiente compatível com a leitura de que a geometria modula o benefício. A Seção~\ref{sec:resposta_qp1} confronta esse gradiente com a família da comparação principal, em que a classificação geométrica não distingue as derrotas.

A Tabela~\ref{tab:hosts_suite} estratifica as mesmas contagens por suíte, a partição em que o Capítulo~\ref{sec:results} encontrou todas as derrotas da família da comparação principal.

\input{generated/tab_hosts_suite}

Nos dois acoplamentos que melhoram o seu hospedeiro, a WFG é a suíte menos favorável. Em IGD, ela tem a menor taxa de vitória tanto do IVF/SPEA2, com 9 vitórias em 18 instâncias, quanto do IVF/NSGA-III, com 3 em 18, e todas as derrotas desses dois acoplamentos, nos dois indicadores, estão nela. O $A_{12}$ mediano acompanha as contagens: em IGD, $0{,}644$ na WFG contra $0{,}769$ nas demais suítes para o IVF/SPEA2, e $0{,}537$ contra $0{,}687$ para o IVF/NSGA-III; em HV, o do IVF/NSGA-III na WFG fica abaixo de $0{,}5$. A exceção à ordenação por taxa de vitória está no HV do IVF/NSGA-III, cuja menor taxa é a da ZDT, com uma vitória em cinco instâncias, seguida da WFG. O IVF/NSGA-II não segue o padrão: em todas as suítes, a maior parte das suas comparações não tem diferença detectável, e é na WFG que ele tem a sua única vitória em IGD e quatro das suas cinco vitórias em HV, mas também uma das suas duas derrotas em IGD e duas das três em HV.

Essa estratificação também é descritiva: nenhum teste de interação entre acoplamento e suíte foi feito. A Seção~\ref{sec:padrao_familias} examina o que a concentração na WFG sustenta quando lida com as demais famílias.

\section{Quando ativar o operador}
\label{sec:ativacao}

As seções anteriores tratam o operador como sempre ativo. Esta reúne a evidência da QP4: é possível decidir, antes ou no início da execução, se vale mantê-lo, e essa decisão compensa diante do IVF/SPEA2 sempre ativo? A análise usa como base os mesmos resultados da comparação principal, o que condiciona a escolha do desfecho.

As instâncias foram rotuladas conforme o operador ajude ou não, a partir da comparação em IGD da coorte da comparação principal: o rótulo registra ajuda quando o teste de Mann--Whitney, sem correção de multiplicidade, rejeita a igualdade a $\alpha = 0{,}05$ e o $A_{12}$ passa de $0{,}56$ a favor do IVF/SPEA2. Avaliar um controlador em IGD reutilizaria os rótulos como sua própria validação. Por isso o HV é o desfecho primário declarado desta família~\cite{zambrano2026clei}. Dos 52 casos disponíveis, um problema de engenharia não recebeu rótulo e ficou fora; as contagens são sobre as \textbf{51 instâncias rotuladas}, 41 como beneficiadas e 10 como não beneficiadas. As 41 excedem as 37 vitórias corrigidas do Capítulo~\ref{sec:results} porque o rótulo dispensa a correção de Holm: WFG1 com $M = 2$ e DTLZ4, DTLZ7 e WFG8 com $M = 3$ são vitórias sem a correção e empates com ela. Entre as dez, seis não exibem diferença detectável e quatro exibem prejuízo detectável em IGD: as quatro derrotas do Capítulo~\ref{sec:results}, as mesmas com e sem correção, todas na WFG.

\subsection{Descritores estáticos de paisagem}

A primeira hipótese testada foi a de que características estáticas da paisagem, medidas antes da execução, permitiriam prever o benefício.

\input{generated/tab_fla_modelos}

A reanálise da Tabela~\ref{tab:ativacao_robustez} aplica aos mesmos 14 descritores uma validação mais estrita: a floresta tem a transformação da resposta ajustada apenas no treino, as predições retornam à escala original e, na retenção por família, DTLZ e MaF saem juntas, o que mantém MaF7 e DTLZ7 do mesmo lado da dobra. A correlação de Spearman agrupada é $0{,}286$ ao reter uma instância e $0{,}150$ quando DTLZ e MaF são retidas juntas. A Tabela~\ref{tab:fla_modelos} preserva os números do estudo original; como a reanálise muda ao mesmo tempo a transformação, o agrupamento das famílias e a escala da resposta, a diferença entre os dois não é atribuída a nenhuma dessas mudanças isoladamente. Nem a regressão reestimada nem os classificadores da tabela original indicam um sinal estável para decidir a ativação fora da família de treino.

\subsection{Dinâmica inicial da população}

A segunda hipótese substitui a medição estática por observáveis colhidos nos primeiros 20\% das gerações.

\input{generated/tab_dinamica_separacao}

A Tabela~\ref{tab:dinamica_separacao} reproduz cinco dos 18 observáveis testados, os de menor valor-$p$; a correção de Benjamini--Hochberg é aplicada sobre os 18. Apenas a renovação do arquivo separa as classes de forma que sobrevive à correção, com $A_{12}$ de $0{,}849$ para a renovação no IVF/SPEA2 e $0{,}824$ para a do SPEA2. As diferenças finais de IGD e de HV entre o acoplamento e o hospedeiro, medidas ao fim das mesmas execuções, não sobrevivem; como só existem ao fim da execução, elas não serviriam de regra de ativação e entram na tabela como referência. O número agregado de ciclos de IVF na fase inicial é semelhante nas duas classes, mas isso não exclui diferenças de intensidade ou de oportunidade por execução.

\input{generated/tab_dinamica_limiar}

\input{generated/tab_ativacao_robustez}

A regra original de limiar único para a renovação do arquivo atinge acurácia balanceada de $0{,}754$ na validação deixa-uma-família por suíte, ante $0{,}691$ em deixa-uma-instância (Tabela~\ref{tab:dinamica_limiar})~\cite{zambrano2026clei}. Em cada dobra, apenas o limiar e sua direção são ajustados sem a família retida; o observável, escolhido entre os sinais examinados, e a janela de 20\% foram fixados na análise de origem sobre esse conjunto, antes das dobras. A reanálise da Tabela~\ref{tab:ativacao_robustez} mantém esse observável e essa janela, retém DTLZ e MaF conjuntamente para impedir a duplicação funcional entre treino e teste e encontra BA $0{,}754$ e MCC $0{,}413$. Com 51 instâncias e apenas dez na classe minoritária, essas estimativas são instáveis, e a concordância não valida de forma independente a escolha do sinal e da janela nem o controlador. Os limiares do protocolo original são $0{,}216$ ao reter DTLZ, MaF ou ZDT e $0{,}248$ ao reter WFG.
A Figura~\ref{fig:ativacao_turnover} mostra a sobreposição entre as distribuições das instâncias rotuladas. Cada ponto é a mediana da renovação inicial nas 30 execuções pareadas por semente de uma instância, e as classes são as 41 instâncias rotuladas como beneficiadas e as 10 não beneficiadas, indicadas na figura por ``Ajuda'' e ``Não ajuda''. A linha tracejada marca $0{,}216$, limiar mediano da validação deixa-uma-família por suíte (Tabela~\ref{tab:dinamica_limiar}); o limiar da WFG é $0{,}248$. A sobreposição das distribuições é compatível com erros de classificação.

\begin{figure}[!htb]
    \centering
    \includegraphics[width=0.85\textwidth]{generated/figures/fig_ativacao_turnover.pdf}
    \caption[Renovação do arquivo na fase inicial, por classe de instância]{Renovação do arquivo na fase inicial, por classe de instância.}
    \label{fig:ativacao_turnover}
\end{figure}

A Tabela~\ref{tab:dinamica_suite} abre a regra por suíte.

\input{generated/tab_dinamica_suite}

Oito das dez instâncias em que o operador não ajuda pertencem à WFG, que tem também a menor renovação inicial mediana das quatro suítes. Aplicados aos \emph{resumos por instância} das 30 execuções, os limiares da validação deixa-uma-família indicam desativação nas 18 instâncias WFG e, fora dela, em DTLZ6 com $M=2$ e $M=3$, ambas rotuladas como beneficiadas. A coincidência com a partição WFG/demais é 49/51 tanto no protocolo original quanto no agrupamento DTLZ+MaF da reanálise; isso impede atribuir a BA $0{,}754$ à discriminação dentro das suítes. Aplicada diretamente aos mesmos rótulos, sem nenhum parâmetro ajustado, a própria partição (desligar na WFG e só nela) atinge BA $0{,}778$ e MCC $0{,}462$ (Tabela~\ref{tab:ativacao_robustez}); formulada depois de observada a concentração das instâncias não beneficiadas, ela serve de referência retrospectiva para a regra de renovação nesta amostra (Seção~\ref{sec:validade_conclusao}). A classificação por resumo usa medianas de trajetórias agrupadas, ao passo que o controlador mede o sinal e decide \emph{em cada execução}, de modo que ela não permite inferir essa decisão. A regra e a partição de suítes acertam as mesmas decisões nesse resumo, sem identificar se a renovação é mecanismo causal ou correlato da suíte.

\subsection{Controlador de aquecimento}

A regra anterior foi implementada como controlador de aquecimento: após 20\% do orçamento, o operador pode ser desativado de forma irreversível conforme a renovação média do arquivo observada nesse período inicial \emph{da própria execução}; os descritores estáticos de paisagem não entram na decisão. O experimento original usou o limiar calibrado por retenção de suíte: $0{,}216$ em DTLZ, MaF e ZDT, e $0{,}248$ na WFG. A reanálise que agrupa DTLZ+MaF não reexecuta nem altera essas decisões do controlador.

\input{generated/tab_controlador_wtl}

Contra o SPEA2 canônico, o controlador registra 34 vitórias, 14 comparações sem diferença detectável e três derrotas em HV nas 51 instâncias. Contra o IVF/SPEA2 sempre ativo, que é o comparador do benefício do controle, o resultado é desfavorável: duas vitórias, 40 comparações sem diferença detectável e nove derrotas em HV.

A decomposição por classe, na mesma tabela, mostra que as nove derrotas diante do IVF/SPEA2 sempre ativo ocorrem em instâncias rotuladas como beneficiadas e as duas vitórias, nas quatro rotuladas como prejudicadas, nas quais o controlador ainda perde para o SPEA2 em três. Nas seis rotuladas como neutras não há diferença detectável contra o operador sempre ativo. Todas as diferenças detectáveis em HV diante do sempre ativo ocorrem na WFG. Como a Tabela~\ref{tab:dinamica_suite} classifica instâncias, nem as suas 18 indicações de desligamento WFG demonstram que o controlador tenha desligado em todas as execuções WFG, nem a indicação em DTLZ6 basta para explicar a derrota observada em IGD com $M=2$.

\section{Posicionamento contra os algoritmos comparados}
\label{sec:posicionamento}

Esta seção reúne a evidência da QP5: como o IVF/SPEA2 se posiciona, em IGD e em HV, diante dos sete algoritmos comparados da avaliação, e em que medida esse posicionamento varia com o número de objetivos. O protocolo é o declarado na Seção~\ref{sec:estatistica}: teste de Mann--Whitney bicaudal por instância, com $\alpha = 0{,}05$ e correção de Holm separadamente por comparador, por número de objetivos e por métrica, com a IGD como desfecho primário e o HV como secundário obrigatório. A coorte é a da comparação principal (o IVF/SPEA2 nas execuções 3001--3060 e cada comparador nas execuções 1--60, 60 execuções por instância e algoritmo), a mesma da Tabela~\ref{tab:confirmatorio_wtl} (Seção~\ref{sec:validade_conclusao}).

A Tabela~\ref{tab:posicionamento} registra as contagens corrigidas, e elas descrevem um posicionamento \textbf{não uniforme}. Chama-se aqui \emph{derrota} a instância com diferença corrigida desfavorável ao IVF/SPEA2, e \emph{saldo desfavorável} o contraste (comparador, número de objetivos e métrica) com mais derrotas do que vitórias. Em IGD com $M = 2$, o saldo favorece o IVF/SPEA2 diante de cinco dos sete algoritmos comparados e é desfavorável diante de dois, o NSGA-III (9/2/17) e o AR-MOEA (9/5/14). Em IGD com $M = 3$, o saldo é favorável diante dos sete, mas estreito diante do AGE-MOEA-II (9/6/8). Em HV com $M = 2$, o saldo é favorável diante dos sete, embora haja derrotas em todos esses contrastes. Em HV com $M = 3$, o saldo é favorável diante de cinco dos sete e desfavorável diante do SPEA2+SDE (8/1/14) e do AGE-MOEA-II (6/1/16).

\input{generated/tab_posicionamento}

A Figura~\ref{fig:posto_medio} resume o mesmo conjunto por outro caminho. Em cada instância, os nove algoritmos são ordenados pela mediana de IGD, e o posto médio agrega essas ordens nas 28 instâncias com $M = 2$ e nas 23 com $M = 3$; a linha sobre cada barra é o desvio padrão do posto entre as instâncias. Com $M = 2$, o IVF/SPEA2 fica em segundo, atrás do NSGA-III, e com $M = 3$, em primeiro. O posto médio é descritivo: ele não considera a magnitude das diferenças nem a variabilidade entre execuções. As contagens corrigidas da tabela e o posto médio da figura chegam a conclusões diferentes porque medem coisas diferentes (frequência de diferença detectável o primeiro, ordenação média o segundo), e é a contagem que sustenta a leitura.

\begin{figure}[!htb]
    \centering
    \includegraphics[width=\textwidth]{generated/figures/fig_posto_medio.pdf}
    \caption[Posto médio exploratório de IGD entre os nove algoritmos]{Posto médio de IGD dos nove algoritmos, por número de objetivos.}
    \label{fig:posto_medio}
\end{figure}

O fecho é por regime de número de objetivos, e não coincide com a ordem média de IGD. Com $M = 2$, os dois saldos desfavoráveis de IGD recaem sobre o NSGA-III e o AR-MOEA, enquanto em HV o saldo é favorável diante dos sete: IGD e HV não concordam nesse regime. Com $M = 3$, a discordância muda de direção. Em IGD o saldo é favorável diante dos sete algoritmos comparados, inclusive dos dois de adaptação de geometria e de referências, com margens desiguais (18/2/3 diante do NSGA-III, 9/6/8 diante do AGE-MOEA-II); em HV, é favorável diante de cinco e desfavorável diante do SPEA2+SDE (8/1/14) e do AGE-MOEA-II (6/1/16). O AGE-MOEA-II e o AR-MOEA, nas configurações padrão da plataforma, não se mostram preferíveis ao IVF/SPEA2 na configuração calibrada: o saldo favorece este em seis dos oito contrastes com eles, com duas exceções, o AR-MOEA em IGD com $M = 2$ e o AGE-MOEA-II em HV com $M = 3$. A leitura restringe-se a esses dois algoritmos assim configurados, sem se estender à classe dos métodos de adaptação de geometria. As contagens de diferença detectável, obtidas em uma coorte e um orçamento específicos, não indicam superioridade em nenhum dos dois regimes, e a discordância entre IGD e HV, que estes contrastes não isolam, fica por explicar.
